\documentclass[a4paper,10pt]{article}

\usepackage[spanish,es-tabla,es-noshorthands]{babel}
\usepackage[margin=2.1cm,top=2.2cm,bottom=2.2cm]{geometry}
\usepackage{amsmath,amssymb}
\usepackage{booktabs,tabularx,array,multirow}
\usepackage{xcolor}
\usepackage{tikz}
\usetikzlibrary{arrows.meta,positioning}
\usepackage{enumitem}
\usepackage{titlesec}
\usepackage{fancyhdr}
\usepackage[
  colorlinks=true,
  linkcolor=blue!45!black,
  citecolor=blue!45!black,
  urlcolor=blue!45!black,
  pdftitle={Comparativa técnica entre HPicoRV32 y HARV},
  pdfauthor={Proyecto HPicoRV32}
]{hyperref}

\definecolor{inkgray}{gray}{0.32}
\definecolor{panelbg}{gray}{0.965}
\definecolor{accentblue}{RGB}{24,64,116}
\definecolor{warnred}{RGB}{150,35,35}

\setlength{\parindent}{0pt}
\setlength{\parskip}{0.42em}
\setlength{\headheight}{13pt}
\setlist{nosep,leftmargin=1.4em}
\setlength{\abovecaptionskip}{4pt}
\setlength{\belowcaptionskip}{4pt}

\titleformat{\section}{\large\bfseries\color{accentblue}}{\thesection.}{0.5em}{}[\vspace{-0.3em}\rule{\textwidth}{0.4pt}]
\titlespacing*{\section}{0pt}{1.0em}{0.4em}
\titleformat{\subsection}{\normalsize\bfseries}{\thesubsection.}{0.5em}{}
\titlespacing*{\subsection}{0pt}{0.7em}{0.25em}

\pagestyle{fancy}
\fancyhf{}
\lhead{\footnotesize\color{inkgray}HPicoRV32 frente a HARV}
\rhead{\footnotesize\color{inkgray}Comparativa técnica de endurecimiento RV32I}
\cfoot{\footnotesize\thepage}
\renewcommand{\headrulewidth}{0.3pt}

\newcolumntype{Y}{>{\raggedright\arraybackslash}X}
\newcolumntype{C}{>{\centering\arraybackslash}X}

\begin{document}

\begin{center}
  {\LARGE\bfseries Comparativa técnica entre HPicoRV32 y HARV}\\[0.25em]
  {\normalsize\color{inkgray} Granularidad TMR, votación y \textit{scrubbing}, ECC en registros/memorias, gestión de \textit{traps} y recursos en Artix-7}\\[0.35em]
  {\footnotesize Proyecto HPicoRV32 \quad$\cdot$\quad Referencia \texttt{external/harv} (\texttt{xarc/harv}) \quad$\cdot$\quad Noviembre de 2026}
\end{center}
\vspace{-0.4em}
\noindent\rule{\textwidth}{0.6pt}

\begin{abstract}
\vspace{-0.4em}
\noindent Este documento compara las técnicas de tolerancia a fallos de \textbf{HPicoRV32} (derivado de PicoRV32 en Verilog-2001) frente al procesador de referencia \textbf{HARV} (VHDL)~\cite{santos2020dtis,harv_repo} sobre FPGA Xilinx Artix-7 (\texttt{xc7a100tcsg324-1}). El contraste abarca la granularidad del TMR, el comportamiento de votadores y \textit{scrubbing}, la codificación ECC de registros y memorias, el manejo de excepciones y las métricas de implementación y de inyección de fallos, explicitando los límites de comparabilidad entre ambos entornos.
\end{abstract}

\section{Arquitectura base y alcance del endurecimiento}

HARV~\cite{santos2020dtis,harv_repo} es un núcleo RV32I escrito en VHDL con ejecución controlada por una máquina de estados de seis fases (\texttt{STAT\_IDLE}, \texttt{STAT\_REQ\_INSTR}, \texttt{STAT\_RUN}, \texttt{STAT\_DMEM\_STALL}, \texttt{STAT\_UPDATE\_PC}, \texttt{STAT\_TRAP}) y camino de datos combinacional encadenado entre \texttt{instr\_fetch}, \texttt{control}, \texttt{regfile}, \texttt{alu} y la actualización del contador de programa (PC). Por su parte, HPicoRV32 parte del núcleo multiciclo PicoRV32~\cite{picorv32} en Verilog-2001 y se integra directamente en un SoC completo con arquitectura Harvard y separación estricta \textit{Write XOR Execute} (\texttt{W\^{}X}): memoria de instrucciones de solo lectura por bus (\texttt{IMEM}, 8~KB), memoria de datos (\texttt{DMEM}, 8~KB), controlador UART y GPIO.

En ambos proyectos la mitigación se configura en tiempo de elaboración mediante parámetros genéricos (\texttt{TMR\_CONTROL}, \texttt{TMR\_ALU}, \texttt{HAMMING\_PC} y \texttt{HAMMING\_REGFILE} en HARV; \texttt{ENABLE\_TMR} y \texttt{ENABLE\_REG\_ECC} en HPicoRV32). La diferencia de control radica en la conmutación en ejecución: HARV expone el registro CSR \texttt{harden\_conf\_r} (\texttt{0xBC0}) para activar o puentear cada técnica desde firmware con el objetivo de medir bajo haz el mismo \textit{bitstream}~\cite{santos2021dtis}. Sin embargo, en \texttt{csr.vhd} este registro carece de cláusula de \textit{reset} (si arranca a cero, la corrección empieza deshabilitada) y no posee redundancia, por lo que un único SEU en \texttt{0xBC0} desactiva el endurecimiento en vuelo. HPicoRV32 fija la protección en síntesis y separa la experimentación mediante el parámetro \texttt{ENABLE\_FI}: el diseño de vuelo (\texttt{FI=0}) elimina toda lógica de inyección con sobrecoste nulo. La Tabla~\ref{tab:arch} resume las diferencias estructurales.

\begin{table}[htbp]
  \centering
  \footnotesize
  \caption{Resumen arquitectónico de ambos diseños.}
  \label{tab:arch}
  \begin{tabularx}{\textwidth}{@{}lYY@{}}
    \toprule
    \textbf{Aspecto} & \textbf{HARV (\texttt{xarc/harv})} & \textbf{HPicoRV32} \\
    \midrule
    Lenguaje y base & VHDL, núcleo propio RV32I (sin sistema) & Verilog-2001, PicoRV32~\cite{picorv32} RV32I \\
    Microarquitectura & Uniciclo/FSM de 6 estados sin segmentar & Multiciclo (8 estados en \texttt{cpu\_state}) \\
    Alcance del proyecto & Núcleo aislado (memorias y SoC externos) & SoC completo (CPU, IMEM, DMEM, UART, GPIO) \\
    Cobertura TMR & Bloques completos \texttt{control} y \texttt{alu} & Grano fino en 5 registros (79~FFs $\times 3$) \\
    Cobertura ECC & PC, IR y \texttt{regfile} (Hamming estándar) & \texttt{cpuregs}, IMEM y DMEM (Hsiao 39,32) \\
    Separación de bus & Puertos Harvard sin barrera \texttt{W\^{}X} interna & Harvard \texttt{W\^{}X} por hardware (\texttt{harv\_err}) \\
    Control dinámico / FI & \textit{Bypass} runtime vía \texttt{0xBC0} (SPOF sin reset) & Fijado en síntesis; bus \texttt{ENABLE\_FI} desacoplado \\
    Portabilidad & Vivado (\texttt{DONT\_TOUCH}) y Synplify/PolarFire & Vivado, \texttt{openXC7} (Yosys) y PolarFire \\
    \bottomrule
  \end{tabularx}
\end{table}

\section{Granularidad del TMR y preservación en síntesis}

HARV aplica una estrategia ``macro'' envolviendo módulos intactos sin modificar su RTL interno~\cite{santos2020dtis}:
\begin{itemize}
  \item \textbf{\texttt{control\_tmr.vhd}:} triplica tres instancias de \texttt{control}, replicando tanto los 3~FFs de estado \texttt{proc\_status\_r} como toda la lógica combinacional de decodificación, y vota sus 24 señales de salida.
  \item \textbf{\texttt{alu\_tmr.vhd}:} triplica tres copias de la \texttt{alu} puramente combinacional y vota bit a bit \texttt{zero\_o} (1~b) y \texttt{data\_o} (32~b). Esta decisión protege frente a transitorios combinacionales (SET) y fallos en LUTs/CRAM, a costa de consumir 1\,559~LUTs solo en \texttt{alu\_tmr} y añadir dos etapas de votación consecutivas dentro del mismo ciclo.
\end{itemize}
Ni PC/IR (cubiertos con Hamming) ni el banco \texttt{csr.vhd} entran en el dominio TMR de HARV.

HPicoRV32 adopta TMR selectivo de grano fino dentro de \texttt{rtl/core/picorv32.v}, triplicando únicamente los cinco registros secuenciales que gobiernan el flujo de control e interfaz de memoria: \texttt{reg\_pc} (32~b), \texttt{reg\_next\_pc} (32~b), \texttt{cpu\_state} (8~b), \texttt{mem\_state} (2~b) y \texttt{latched\_rd} (5~b), es decir, $79 \times 3 = 237$~FFs. La ALU y la lógica combinacional no se triplican. Esta elección recorta drásticamente el consumo de LUTs y permite cumplir los 100~MHz en el SoC, asumiendo como compromiso explícito que un pulso SET en la ALU combinacional no queda enmascarado si coincide con el flanco de captura.

Para impedir que el sintetizador fusione réplicas idénticas, HARV anota cada instancia con el atributo \texttt{DONT\_TOUCH} (Vivado) y \texttt{syn\_radhardlevel="tmr"}, \texttt{syn\_keep}, \texttt{syn\_noprune} y \texttt{syn\_safe\_case} (Synplify para Microchip PolarFire). HPicoRV32 aplica atributos equivalentes a nivel de registro (\texttt{dont\_touch="true"}, \texttt{syn\_preserve=1}, \texttt{keep="true"}) y en el módulo votador (\texttt{tmr\_voter.v}: \texttt{keep}, \texttt{syn\_keep}). Además, dado que Yosys no propaga el atributo \texttt{keep} desde declaraciones \texttt{reg} hacia las celdas \texttt{DFF} tras \texttt{proc\_dff}, el flujo abierto \texttt{openXC7} ejecuta \texttt{setattr -set keep 1 a:keep \%ci1 t:*dff* \%i} antes de \texttt{synth\_xilinx}, verificando en \texttt{soc.json} las tres copias físicas exactas (\texttt{96/96/24/6/15} bits sin fusión).

\section{Topología de votadores y \textit{scrubbing} (\textit{self-healing})}

En la salida, ambos diseños utilizan votadores mayoritarios combinacionales bit a bit, $(r_0 \land r_1) \lor (r_0 \land r_2) \lor (r_1 \land r_2)$, sin triplicar el propio votador y acompañados de una red XOR simétrica que alimenta los contadores de discrepancias. En HARV, cada salida votada atraviesa además un multiplexor 2:1 controlado por \texttt{correct\_error\_i} (\texttt{harden\_conf\_r[4..5]}), añadiendo un nivel lógico adicional tras el votador. En \texttt{tmr\_voter.v} de HPicoRV32 el valor votado sale directamente sin multiplexor de \textit{bypass}.

La diferencia funcional crítica reside en la realimentación del estado votado en máquinas secuenciales (Figura~\ref{fig:tmr_sh}):
\begin{itemize}
  \item \textbf{HARV sin \textit{self-healing}:} en \texttt{control\_tmr.vhd}, cada réplica calcula \texttt{next\_proc\_status\_w(i)} a partir de su propio registro interno \texttt{proc\_status\_r(i)}. Ante un SEU en la réplica~0, las réplicas 1 y 2 enmascaran el fallo hacia el resto del procesador, pero la réplica~0 nunca recupera el estado votado: diverge de forma permanente (o queda atascada en \texttt{STAT\_TRAP}) hasta un \textit{reset} global, incrementa de continuo el contador \texttt{0x7C8} y deja al sistema sin redundancia frente a un segundo SEU en otra copia.
  \item \textbf{HPicoRV32 con \textit{self-healing} en 1 ciclo:} dentro de \texttt{always @(posedge clk)} en \texttt{picorv32.v}, cada ciclo precarga \texttt{next\_<reg> = <voted>} con la salida mayoritaria y actualiza simultáneamente las tres réplicas (\texttt{\_r0, \_r1, \_r2 <= next\_<reg>}). Cualquier discrepancia en un biestable se corrige automáticamente en el siguiente flanco de reloj sin ciclos de parada.
\end{itemize}

\begin{figure}[htbp]
  \centering
  \begin{tikzpicture}[
    >= {Stealth[length=1.8mm,width=1.25mm]},
    box/.style={draw=black,line width=0.65pt,rounded corners=2.5pt,align=center,font=\scriptsize,fill=white,inner sep=2.5pt},
    panel/.style={draw=black!55,line width=0.65pt,rounded corners=4pt,fill=panelbg},
    wire/.style={->,draw=black,line width=0.65pt},
    fb/.style={->,draw=accentblue,line width=0.95pt}
  ]
    % Panel (a): HARV (0.0 .. 7.8)
    \draw[panel] (0.0,0.0) rectangle (7.8,5.25);
    \node[font=\footnotesize\bfseries] at (3.90,4.95) {(a) HARV: \texttt{control\_tmr.vhd} (sin \textit{self-healing})};

    \node[box,minimum width=1.95cm,minimum height=0.54cm] (ha0) at (1.78,3.68) {\texttt{proc\_status\_r(0)}};
    \node[box,minimum width=1.95cm,minimum height=0.54cm] (ha1) at (1.78,2.75) {\texttt{proc\_status\_r(1)}};
    \node[box,minimum width=1.95cm,minimum height=0.54cm] (ha2) at (1.78,1.82) {\texttt{proc\_status\_r(2)}};

    \foreach \n/\y in {ha0/3.68, ha1/2.75, ha2/1.82} {
      \draw[wire,draw=warnred] (2.90,\y) -- (2.90,{\y+0.39}) -- (0.50,{\y+0.39}) -- (0.50,\y) -- (\n.west);
    }
    \node[font=\fontsize{6.0}{7.0}\selectfont,text=warnred] at (1.80,4.34) {bucle local: \texttt{next(i) <- r(i)}};

    \draw (ha0.east) -- (3.14,3.68);
    \draw (ha1.east) -- (3.14,2.75);
    \draw (ha2.east) -- (3.14,1.82);
    \draw (3.14,1.50) -- (3.14,3.68);
    \fill (3.14,2.75) circle (1.1pt);
    \fill (3.14,3.15) circle (1.1pt);

    \node[box,minimum width=1.68cm,minimum height=1.08cm] (va) at (4.42,3.15) {\textbf{Votador}\\2-de-3\\(24 salidas)};
    \node[box,minimum width=1.68cm,minimum height=0.70cm] (xa) at (4.42,1.50) {\textbf{Red XOR}\\CSR \texttt{0x7C8}};
    \node[box,minimum width=1.30cm,minimum height=1.02cm] (ma) at (6.36,3.15) {\textbf{Mux 2:1}\\\textit{bypass}};
    \node[box,minimum width=1.65cm,minimum height=0.70cm,draw=warnred] (ca) at (6.36,1.50) {\texttt{harden\_conf[4]}\\{\fontsize{5.8}{6.8}\selectfont\color{warnred}SPOF sin reset}};

    \draw[wire] (3.14,3.15) -- (va.west);
    \draw[wire] (3.14,1.50) -- (xa.west);
    \draw[wire] (va.east) -- (ma.west);
    \draw[wire,draw=warnred] (ca.north) -- (ma.south);
    \draw[wire] (ma.east) -- (7.64,3.15);
    \node[font=\fontsize{6.0}{7.0}\selectfont] at (7.36,3.38) {out};
    \node[font=\fontsize{6.0}{7.0}\selectfont\itshape,text=warnred] at (3.90,0.45) {Un SEU en \texttt{\_r(i)} desincroniza la réplica hasta reset global.};

    % Panel (b): HPicoRV32 (8.3 .. 16.1)
    \draw[panel] (8.3,0.0) rectangle (16.1,5.25);
    \node[font=\footnotesize\bfseries] at (12.20,4.95) {(b) HPicoRV32: TMR + \textit{self-healing} (1 ciclo)};

    \node[box,minimum width=1.95cm,minimum height=0.54cm] (hb0) at (10.42,3.68) {\texttt{<reg>\_r0}};
    \node[box,minimum width=1.95cm,minimum height=0.54cm] (hb1) at (10.42,2.75) {\texttt{<reg>\_r1}};
    \node[box,minimum width=1.95cm,minimum height=0.54cm] (hb2) at (10.42,1.82) {\texttt{<reg>\_r2}};

    \draw (hb0.east) -- (11.78,3.68);
    \draw (hb1.east) -- (11.78,2.75);
    \draw (hb2.east) -- (11.78,1.82);
    \draw (11.78,1.50) -- (11.78,3.68);
    \fill (11.78,2.75) circle (1.1pt);
    \fill (11.78,3.15) circle (1.1pt);

    \node[box,minimum width=1.82cm,minimum height=1.08cm] (vb) at (13.15,3.15) {\textbf{\texttt{tmr\_voter}}\\2-de-3\\(sin \textit{bypass})};
    \node[box,minimum width=1.82cm,minimum height=0.70cm] (xb) at (13.15,1.50) {\textbf{Red XOR}\\CSR \texttt{0xCC0/1}};

    \draw[wire] (11.78,3.15) -- (vb.west);
    \draw[wire] (11.78,1.50) -- (xb.west);
    \draw[wire] (vb.east) -- (15.92,3.15);
    \node[font=\fontsize{6.0}{7.0}\selectfont] at (15.36,3.38) {\texttt{voted}};

    % Lazo self-healing
    \fill[accentblue] (14.62,3.15) circle (1.3pt);
    \draw[line width=0.95pt,draw=accentblue] (14.62,3.15) -- (14.62,4.34) -- (8.80,4.34) -- (8.80,1.82);
    \draw[fb] (8.80,3.68) -- (hb0.west);
    \draw[fb] (8.80,2.75) -- (hb1.west);
    \draw[fb] (8.80,1.82) -- (hb2.west);
    \node[font=\fontsize{6.0}{7.0}\selectfont\bfseries,text=accentblue] at (11.75,4.54) {\texttt{next\_<reg> = voted} (precarga síncrona)};
    \node[font=\fontsize{6.0}{7.0}\selectfont\itshape,text=accentblue] at (12.20,0.45) {Captura simultánea \texttt{\_r0..\_r2 <= next\_<reg>}: repara en 1 ciclo.};
  \end{tikzpicture}
  \caption{Comparación estructural del lazo de realimentación TMR entre HARV y HPicoRV32.}
  \label{fig:tmr_sh}
\end{figure}

En el dominio ECC ocurre el mismo contraste. HARV decodifica y corrige únicamente de forma combinacional en la salida de \texttt{hamming\_register} y \texttt{regfile}, sin reescribir la celda dañada; en registros de vida larga un segundo SEU convierte el fallo en doble error incorregible (DBE). HPicoRV32 ejecuta \textit{auto-scrubbing} automático en todos sus bancos:
\begin{enumerate}
  \item \textbf{\texttt{cpuregs} (banco de registros):} cuando la CPU no realiza escritura (\texttt{!cpuregs\_cpu\_wen}) y algún puerto de lectura presenta \texttt{sbiterr}, reescribe la palabra completa de 39~bits reparada en el mismo flanco de reloj (\texttt{cpuregs\_scrub1} con prioridad sobre \texttt{cpuregs\_scrub2}), con cero ciclos de penalización.
  \item \textbf{BRAM (\texttt{bram\_ecc.v}):} ante un error simple (\texttt{SBE}) en lectura (\texttt{ST\_READ}), transita a \texttt{ST\_SCRUB} (3 ciclos totales frente a 2 en lectura limpia) y reescribe \texttt{scrub\_word\_q} en la memoria física, incluso en \texttt{IMEM} pese a tener \texttt{READ\_ONLY=1} en el bus.
\end{enumerate}

\section{Protección de registros y memorias ECC}

Ambos procesadores adoptan códigos SEC-DED (39,32) sistemáticos (32~bits de dato + 7~bits de redundancia), pero la construcción de la matriz de paridad y la arquitectura de almacenamiento determinan su coste y fiabilidad en FPGA (Tabla~\ref{tab:ecc}).

\begin{table}[htbp]
  \centering
  \footnotesize
  \caption{Comparativa del subsistema ECC SEC-DED (39,32).}
  \label{tab:ecc}
  \begin{tabularx}{\textwidth}{@{}lYY@{}}
    \toprule
    \textbf{Característica} & \textbf{HARV} & \textbf{HPicoRV32} \\
    \midrule
    Código SEC-DED & Hamming clásico + paridad global \texttt{[38]} & Matriz Hsiao (39,32)~\cite{hsiao1970} balanceada \\
    Ecuaciones XOR & Desbalanceadas (6--18 bits) + \texttt{xor\_reduce(39b)} & 13--14 entradas por fila, columnas peso 3 \\
    \textit{Auto-scrubbing} & Ausente (solo corrección combinacional) & En \texttt{cpuregs} (1c), IMEM y DMEM (\texttt{ST\_SCRUB}) \\
    Banco de registros & 1\,209~FFs + 306 \texttt{MUXF7} (\texttt{31$\times$39}) & 14 \texttt{RAM32M} en LUTRAM (\texttt{32$\times$39}) \\
    Registros PC / IR & \texttt{hamming\_register} (IR con flags \texttt{open}) & PC con TMR; IR protegido vía IMEM ECC \\
    Memorias IMEM / DMEM & Externas al núcleo (sin ECC ni RMW en RTL) & 8~KB + 8~KB en \texttt{bram\_ecc.v} con RMW 3c (\texttt{sb/sh}) \\
    Falsos positivos & Lee \texttt{rs1/rs2} aunque no se usen; bit \texttt{[38]} ciego & Cualifica \texttt{reg\_uses\_rs1/2}; paridad simétrica \\
    \bottomrule
  \end{tabularx}
\end{table}

El análisis directo del RTL de HARV (\texttt{external/harv/hdl/ft\_components/}) revela cuatro puntos débiles relevantes frente al diseño de HPicoRV32:
\begin{enumerate}
  \item \textbf{Matriz desbalanceada y bit global \texttt{[38]}:} las filas de \texttt{PARITY\_POS} en \texttt{hamming\_pkg.vhd} agrupan entre 6 y 18~datos y requieren un árbol \texttt{xor\_reduce} de 38/39 bits para la paridad total. Además, un bit-\textit{flip} en \texttt{[38]} deja el síndrome \texttt{parity\_diff\_w} a cero, manteniendo \texttt{single\_err\_o=0} sin notificar hasta que un segundo SEU provoca un fallo doble directo. En HPicoRV32 (\texttt{rtl/core/ecc\_secded.v}), la matriz Hsiao (39,32)~\cite{hsiao1970} reparte uniformemente 13--14 entradas por ecuación, trata por igual los 39~bits (corrigiendo también \texttt{ecc\_o[6:0]} para \textit{scrubbing}) y garantiza que \texttt{32'h0} codifique a \texttt{39'h0} tras el arranque en FPGA.
  \item \textbf{\textit{Bug} de reset no nulo (\texttt{hamming\_pkg.vhd:123--125}):} la función auxiliar \texttt{get\_parity\_data\_result} devuelve \texttt{xor\_reduce(par\_data)} concatenado con \texttt{par\_data} sin incluir \texttt{data} en el bit~38. Con dirección de arranque nula queda oculto, pero cualquier \texttt{RESET\_VALUE} con peso impar inicializa invertida la paridad global y convierte el primer SEU en un falso error doble sin corrección.
  \item \textbf{Inferencia de \texttt{regfile} en Artix-7:} en HARV, el banco \texttt{g\_HAMMING} fuerza la síntesis en 1\,209 biestables individuales con árboles \texttt{MUXF7} para lectura combinacional (y en \texttt{g\_NORMAL} la comparación metalógica \texttt{rd\_i /= "UUUUU"} provoca que Vivado pode el banco completo al sintetizar el core aislado). HPicoRV32 conserva un único puerto de escritura síncrono y dos lecturas asíncronas en \texttt{cpuregs}, manteniendo la inferencia de memoria distribuida (\texttt{14~RAM32M}, solo $+2$ frente a los 12 del diseño sin ECC).
  \item \textbf{Lecturas especulativas y cobertura de memoria:} HARV decodifica siempre \texttt{rs1} y \texttt{rs2} aunque la instrucción (p.~ej., \texttt{LUI} o \texttt{ADDI}) ignore esos campos, generando falsas alarmas; HPicoRV32 filtra \texttt{reg\_dbe\_event} con \texttt{reg\_uses\_rs1/2} (excluyendo también lecturas CSR y \texttt{mret}). Asimismo, HARV delega el ECC de \texttt{dmem} al SoC externo, deja sin protección ni telemetría a \texttt{imem} y desconecta (\texttt{open}) las señales de error del registro de instrucción (\texttt{register\_instr\_u}). HPicoRV32 integra \texttt{bram\_ecc.v} para IMEM y DMEM: registra \texttt{rdata\_o} para cortar el camino crítico desde BRAM, ejecuta \textit{Read-Modify-Write} de 3 ciclos para \texttt{sb/sh} (saneando SBE previo y abortando la escritura si hay DBE), soporta backend dual (RTL Hsiao portable o IP \texttt{blk\_mem\_gen} \textit{Soft ECC} en Vivado mediante \texttt{VIVADO\_ECC\_IP=1}) y aplica barrera \texttt{W\^{}X}.
\end{enumerate}

\section{Gestión de \textit{traps} y telemetría}

HARV dispone de un único mecanismo de \textit{trap} en \texttt{control.vhd} (\texttt{STAT\_TRAP} $\rightarrow$ \texttt{TRAP\_HANDLER\_ADDR}), capturando \texttt{uscratch}, \texttt{uepc}, \texttt{ucause}, \texttt{utval} y \texttt{uip} en \texttt{csr.vhd}. Su revisión muestra limitaciones importantes: en \texttt{control.vhd:172} la condición \texttt{or dmem\_sbu\_i = '1' or dmem\_dbu\_i = '1'} está comentada, de modo que los códigos \texttt{ucause} \texttt{0x10} y \texttt{0x20} nunca disparan \textit{trap}; \texttt{ecall\_o} está fijado a \texttt{'0'}; no detecta instrucciones ilegales; \texttt{imem\_err\_i} bloquea la máquina en \texttt{STAT\_IDLE}; no implementa instrucción de retorno (\texttt{mret/uret}) ni guarda frente a reentrancia. Los 10 contadores de error (\texttt{0x7C0}--\texttt{0x7C9}) y los registros de excepción son de lectura/escritura y carecen de protección.

HPicoRV32 separa explícitamente un mecanismo de \textbf{Soft-Trap} (recuperable o diagnosticable por software) de un \textbf{Hard-Trap} (parada segura inmediata en \texttt{cpu\_state\_trap} con \texttt{LED15}), evitando activar \texttt{ENABLE\_IRQ/QREGS} en PicoRV32 porque ampliar \texttt{cpuregs} de 32 a 36 palabras duplicaría el uso de LUTRAM de 14 a 28~\texttt{RAM32M} (Tabla~\ref{tab:traps}).

\begin{table}[htbp]
  \centering
  \footnotesize
  \caption{Mecanismos de excepción y telemetría de fallos.}
  \label{tab:traps}
  \begin{tabularx}{\textwidth}{@{}lYY@{}}
    \toprule
    \textbf{Dimensión} & \textbf{HARV} & \textbf{HPicoRV32} \\
    \midrule
    Niveles de excepción & Único (\texttt{STAT\_TRAP} $\rightarrow$ \texttt{TRAP\_HANDLER\_ADDR}) & Dual: Soft-Trap (\texttt{0x0000\_0010}) + Hard-Trap \\
    Causas capturadas & Solo fallo de bus en \texttt{dmem} (\texttt{dmem\_sbu/dbu} comentados; \texttt{ecall=0}) & \texttt{ILLINSN} (1), \texttt{BREAK} (2), \texttt{MISALIGN} (3), \texttt{HARV\_WX} (4), \texttt{DMEM\_DBE} (5) \\
    DBE en \texttt{regfile} o \texttt{imem} & Solo incrementa contador (no detiene ni salta) & Hard-Trap inmediato (\texttt{cpu\_trap=1}), bloquea DMEM/LEDs \\
    Retorno y reentrancia & Sin \texttt{mret}/\texttt{uret} ni bloqueo de reentrancia & \texttt{mret} (\texttt{0x3020\_0073}); segunda excepción escala a Hard-Trap \\
    Telemetría de errores & 10 CSRs R/W (\texttt{0x7C0}--\texttt{0x7C9}, incl.\ ALU y \texttt{dmem}) & 4 CSRs R/O (\texttt{0xCC0}--\texttt{0xCC3}) + 4 MMIO R/O (\texttt{0x3000\_0000..0C}) \\
    Diagnóstico \textit{trap} & \texttt{uepc}/\texttt{ucause}/\texttt{utval} en \texttt{csr.vhd} (sin proteger) & \texttt{CSR\_TRAP\_EPC} (\texttt{0xCC4}) y \texttt{CSR\_TRAP\_CAUSE} (\texttt{0xCC5}) R/O \\
    \bottomrule
  \end{tabularx}
\end{table}

Cuando ocurre una de las cinco causas capturables con \texttt{trap\_active=0} y sin fallos fatales en \texttt{cpuregs} ni \texttt{IMEM}, HPicoRV32 graba \texttt{trap\_epc} (\texttt{0xCC4}) y \texttt{trap\_cause} (\texttt{0xCC5}), activa \texttt{trap\_active}, cancela transacciones en curso y salta a \texttt{0x0000\_0010} usando el camino del PC protegido por TMR. El vector \texttt{\_trap\_vector} en \texttt{crt0.S} preserva los 16 registros \textit{caller-saved} en \texttt{DMEM} (habilitada mientras \texttt{cpu\_trap=0}), invoca a \texttt{trap\_handler()} en C para reportar causa, dirección y los ocho contadores por LEDs y UART, y retorna con \texttt{mret} ante \texttt{BREAK} o detiene el sistema. Si se produce un DBE en \texttt{cpuregs} o \texttt{IMEM} —donde no es seguro ejecutar código ni confiar en la pila— o una segunda excepción anidada (\texttt{trap\_active=1}), la CPU entra en Hard-Trap sin sobrescribir los registros de diagnóstico.

\section{Consumo de recursos, temporización y verificación}

La Tabla~\ref{tab:res} reúne las métricas \textit{post-route} sobre Artix-7 \texttt{xc7a100tcsg324-1} con reloj objetivo de 100~MHz. Para interpretarlas con rigor deben tenerse en cuenta dos matices metodológicos: (1) las corridas de HARV en Vivado~2025.2 corresponden a síntesis \textit{Out-Of-Context} (OOC) del núcleo aislado (sin SoC ni memorias y con \texttt{regfile} podado en \texttt{harv\_base}, \texttt{harv\_tmr\_ctrl} y \texttt{harv\_tmr\_all}), mientras que las de HPicoRV32 incluyen el SoC completo; (2) las cifras de Vivado para HPicoRV32 corresponden al Hito~2 (pre-\textit{Soft-Trap}), mientras que la versión final consolidada con \textit{Soft-Trap} se midió en el flujo abierto \texttt{openXC7} (\texttt{Yosys 0.69} + \texttt{nextpnr-xilinx}).

\begin{table}[htbp]
  \centering
  \footnotesize
  \setlength{\tabcolsep}{4pt}
  \caption{Recursos y temporización \textit{post-route} en Artix-7 (\texttt{xc7a100tcsg324-1} @ 100~MHz).}
  \label{tab:res}
  \begin{tabularx}{\textwidth}{@{}Ylrrrrrr@{}}
    \toprule
    \textbf{Diseño / Flujo} & \textbf{Ámbito} & \textbf{LUTs} & \textbf{FFs} & \textbf{LUTRAM} & \textbf{BRAM} & \textbf{WNS (ns)} & \textbf{$F_{\max}$ (MHz)} \\
    \midrule
    \texttt{harv\_base} (Vivado)$^\dagger$ & Core OOC & 845 & 136 & 0 & 0 & $-3.141$ & $\sim 76.1$ \\
    \texttt{harv\_tmr\_ctrl} (Vivado)$^\dagger$ & Core OOC & 928 & 147 & 0 & 0 & $-4.322$ & $\sim 69.8$ \\
    \texttt{harv\_tmr\_all} (Vivado)$^\dagger$ & Core OOC & 1\,995 & 195 & 0 & 0 & $-5.515$ & $\sim 64.5$ \\
    \texttt{harv\_ecc} (Vivado) & Core OOC & 2\,674 & 1\,808 & 0 & 0 & $-14.964$ & $\sim 40.1$ \\
    \texttt{harv\_ft\_full} (Vivado) & Core OOC & 3\,958 & 1\,814 & 0 & 0 & $-17.504$ & $\sim 36.4$ \\
    \midrule
    HPicoRV32 Base (Vivado H0) & SoC & 1\,114 & 703 & 44 & 4 & $+3.531$ & $\sim 154.6$ \\
    HPicoRV32 TMR (Vivado H1) & SoC & 1\,429 & 936 & 44 & 4 & $+2.298$ & $\sim 129.8$ \\
    HPicoRV32 TMR+ECC RTL (Vivado H2) & SoC & 2\,135 & 1\,375 & 56 & 5 & $+0.226$ & $\sim 102.3$ \\
    HPicoRV32 TMR+ECC IP (Vivado H2) & SoC & 2\,111 & 1\,368 & 56 & 5 & $+0.240$ & $\sim 102.5$ \\
    \midrule
    HPicoRV32 Soft-Trap (\texttt{openXC7}, \texttt{FI=0}) & SoC & 2\,684$^\ddagger$ & 1\,401 & 14~\texttt{RAM32M} & 5 eq. & $+0.778$ & $108.44$ \\
    \bottomrule
  \end{tabularx}
  \vspace{3pt}
  \begin{minipage}{\textwidth}
    \raggedright\scriptsize
    $^\dagger$En estas tres variantes OOC, Vivado podó \texttt{regfile\_u} (\texttt{g\_NORMAL} incluye \texttt{rd\_i /= "UUUUU"}).
    $^\ddagger$LUTs lógicas de Yosys (\texttt{4\,055 SLICE\_LUTX} en \texttt{nextpnr}). Con \texttt{FI=1} (\texttt{soc\_fi\_top}) consume 1\,527~FFs y cierra a 103.49~MHz (\textit{slack} $+337$~ps).
  \end{minipage}
\end{table}

Pese a incluir memorias ECC, UART y decodificación \texttt{W\^{}X}, el SoC HPicoRV32 completo en Vivado (\texttt{2\,135} LUTs y \texttt{1\,375} FFs) requiere un \textbf{46\,\% menos de LUTs} y un \textbf{24\,\% menos de FFs} que el núcleo aislado \texttt{harv\_ft\_full} (\texttt{3\,958} LUTs y \texttt{1\,814} FFs), operando casi \textbf{2.8 veces más rápido} ($102.3$--$108.4$~MHz frente a $36.4$~MHz). Tres decisiones explican esta diferencia: (1) la ejecución multiciclo evita encadenar en un único periodo la decodificación del IR, lectura del banco, ALU y recodificación (\texttt{27.380~ns}, con 82.3\,\% de retardo de enrutado en \texttt{harv\_ft\_full}); (2) el uso de LUTRAM en \texttt{cpuregs} elimina 1\,209~FFs y más de 300~\texttt{MUXF7}; (3) omitir \texttt{alu\_tmr} ahorra $\sim$1\,559~LUTs y nueve niveles lógicos extra, mientras el registro \texttt{rdata\_o} de \texttt{bram\_ecc} aísla la ruta crítica de las BRAM (que solo crecen de 4 a 5 tiles equivalentes). El soporte completo de \textit{Soft-Trap} en \texttt{openXC7} añade únicamente $+67$~LUTs Yosys y $+41$~FFs sin afectar a LUTRAM ni BRAM.

Para contrastar la eficacia frente a fallos se ejecutó el mismo \textit{benchmark} de suma de vectores de 300 elementos usado en la evaluación SoC de HARV~\cite{santos2021microrel} ($N=1\,000$ experimentos por escenario). De nuevo es necesario distinguir el alcance de la inyección: el artículo de HARV inyectó SEUs mediante ModelSim sobre todos los biestables del procesador y periféricos en Zynq (1\,453~FFs en el SoC endurecido, excluyendo IMEM/DMEM), obteniendo una subida de \textit{Match} del \textbf{90.9\,\% al 99.3\,\%} en 1-SEU y del \textbf{51.8\,\% al 93.9\,\%} en 10-SEU, con fallos residuales en el bus y la UART sin proteger. En HPicoRV32 (\texttt{tb\_fault\_campaign.v} en Verilator), la campaña de Fase~1 inyecta mediante el bus hardware desacoplado sobre los 8 objetivos endurecidos (perfiles CPU \texttt{0..5} y SoC \texttt{0..7} incluyendo IMEM/DMEM): la línea base sin TMR ni ECC en registros (\texttt{0,0}) obtiene \textbf{29.7\,\%} de \textit{Match} en 1-SEU y \textbf{0.0\,\%} en 10-SEU (99.5\,\% \textit{Hang}), mientras que la configuración completa (\texttt{1,1}) alcanza \textbf{100.0\,\% de \textit{Match}} (\texttt{0} SDC, \texttt{0} \textit{Hang}) tanto en 1-SEU como en 10-SEU sobre los perfiles CPU y SoC, además de 4/4 \textit{traps} controlados ante DBE dirigidos en \texttt{cpuregs}, IMEM, DMEM y violación \texttt{W\^{}X}. Queda pendiente para una Fase~2 extender la inyección externa de Verilator al resto de biestables no protegidos del SoC (bus MMIO y UART) para obtener la cifra global equivalente a~\cite{santos2021microrel}.

\section{Síntesis de compromisos y frentes abiertos}

HPicoRV32 corrige las principales carencias estructurales observadas en el código de referencia de HARV —ausencia de \textit{self-healing} en \texttt{control\_tmr}, falta de \textit{scrubbing} ECC, banco de registros en FFs, ruta combinacional uniciclo de 36~MHz, \textit{bug} de paridad global en reset, memoria sin ECC/\texttt{W\^{}X} y SPOF en \texttt{0xBC0}— logrando cerrar temporización por encima de 100~MHz con menos área en FPGA. A cambio, asume tres compromisos explícitos que conviene vigilar antes de un despliegue final en misión:
\begin{itemize}
  \item \textbf{Exposición de la ALU a SET:} al no triplicar la lógica combinacional de cálculo para mantener área y frecuencia, la protección frente a transitorios en LUTs es menor que en \texttt{harv\_tmr\_all}.
  \item \textbf{Registros de telemetría y periféricos sin endurecer:} igual que \texttt{csr.vhd} en HARV, los contadores CSR/MMIO, el periférico UART y los registros auxiliares del camino de datos (\texttt{reg\_op1}, \texttt{reg\_op2}, \texttt{reg\_out}) permanecen en FFs simples.
  \item \textbf{Cobertura total del SoC en inyección:} la Fase~1 valida los objetivos endurecidos a 100~MHz en simulación y placa, pero requiere la Fase~2 sobre el 100\,\% de los biestables del SoC para cuantificar la sección eficaz residual de los bloques sin proteger.
\end{itemize}

\bibliographystyle{unsrt}
\small
\bibliography{references}

\end{document}
